\documentclass[12pt, a4paper]{article}

% =========================================================================
% PACOTES E CONFIGURAÇÕES
% =========================================================================
\usepackage[utf8]{inputenc}
\usepackage[T1]{fontenc}
\usepackage[brazil]{babel}
\usepackage{geometry}
\geometry{a4paper, left=3cm, top=3cm, right=2cm, bottom=2cm}
\usepackage{setspace}
\onehalfspacing
\usepackage{indentfirst}
\usepackage{graphicx}
\usepackage{xcolor}
\usepackage{hyperref}
\usepackage{listings}
\usepackage{booktabs}
\usepackage{amsmath}
\usepackage{tikz}
\usetikzlibrary{shapes.geometric, arrows.meta, positioning, fit, backgrounds, calc}

% Configuração de links
\hypersetup{
    colorlinks=true,
    linkcolor=blue!70!black,
    urlcolor=blue!70!black,
    citecolor=blue!70!black
}

% Configuração de blocos de código
\definecolor{codebg}{rgb}{0.96,0.96,0.96}
\definecolor{codeframe}{rgb}{0.8,0.8,0.8}
\definecolor{codegreen}{rgb}{0,0.5,0}
\definecolor{codeblue}{rgb}{0.1,0.2,0.7}

\lstset{
    backgroundcolor=\color{codebg},
    frame=single,
    rulecolor=\color{codeframe},
    basicstyle=\ttfamily\footnotesize,
    breaklines=true,
    keywordstyle=\color{codeblue}\bfseries,
    commentstyle=\color{codegreen}\itshape,
    stringstyle=\color{red!70!black},
    showstringspaces=false,
    tabsize=4,
    inputencoding=utf8,
    extendedchars=true,
    literate={á}{{\'a}}1 {ã}{{\~a}}1 {â}{{\^a}}1 {à}{{\`a}}1
             {é}{{\'e}}1 {ê}{{\^e}}1 {í}{{\'i}}1 {ó}{{\'o}}1
             {õ}{{\~o}}1 {ô}{{\^o}}1 {ú}{{\'u}}1 {ç}{{\c{c}}}1
             {Á}{{\'A}}1 {Ã}{{\~A}}1 {É}{{\'E}}1 {Í}{{\'I}}1
             {Ó}{{\'O}}1 {Õ}{{\~O}}1 {Ú}{{\'U}}1 {Ç}{{\c{C}}}1
}

% =========================================================================
% INÍCIO DO DOCUMENTO
% =========================================================================
\begin{document}

% -------------------------------------------------------------------------
% CAPA / CABEÇALHO
% -------------------------------------------------------------------------
\begin{titlepage}
    \centering
    \vspace*{1cm}
    {\Large \textbf{DISCIPLINA DE SISTEMAS DISTRIBUÍDOS}}\\[1.5cm]
    
    {\huge \textbf{Sistema de Armazenamento Distribuído com Fragmentação Binária, Exclusão Mútua via Token Ring, Deduplicação SHA-256 e Tolerância a Falhas}}\\[1.5cm]
    
    \textbf{Relatório Técnico do Projeto Prático}\\[2cm]
    
    \begin{flushright}
        \textbf{Integrantes do Grupo:}\\
        Guilherme e Equipe\\[0.5cm]
        \textbf{Tecnologias Utilizadas:}\\
        Python 3.11, FastAPI (REST + Swagger), gRPC (HTTP/2 + Protobuf),\\
        PyJWT (Stateless Auth), SQLite, Multithreading, SHA-256 e Docker Compose.
    \end{flushright}
    
    \vfill
    {\large \today}
\end{titlepage}

\tableofcontents
\newpage

% =========================================================================
% 1. INTRODUÇÃO
% =========================================================================
\section{Introdução}

O crescimento exponencial do volume de dados trafegados em aplicações modernas tornou inviável o armazenamento e o processamento de arquivos em arquiteturas monolíticas e centralizadas. Sistemas de arquivos distribuídos de larga escala --- como o \textit{Google File System} (GFS) e o \textit{Hadoop Distributed File System} (HDFS) --- demonstram que a divisão de arquivos em fragmentos (\textit{chunks}), associada à distribuição de carga, controle rigoroso de concorrência e replicação de dados, é fundamental para alcançar alta disponibilidade, escalabilidade e integridade.

Neste trabalho, foi projetado e desenvolvido um \textbf{Sistema Distribuído Completo de Armazenamento de Arquivos}, conteinerizado através do \textit{Docker Compose} e dividido em quatro nós independentes comunicando-se em rede: um \textbf{Servidor de Borda (\textit{Edge Server})}, dois \textbf{Servidores Centrais de Processamento e Armazenamento (\textit{Central-1} e \textit{Central-2})} e um \textbf{Servidor de Redundância (\textit{Server-Backup})}.

A solução integra conceitos fundamentais abordados na disciplina de Sistemas Distribuídos:
\begin{itemize}
    \item \textbf{Arquitetura \textit{Stateless} e Tradução de Protocolos:} Exposição de uma API REST (HTTP/JSON e Multipart) com autenticação \textit{Stateless} via \textbf{JWT (\textit{JSON Web Token})} na borda, convertendo as chamadas para \textbf{gRPC (HTTP/2 e \textit{Protocol Buffers})} na rede interna;
    \item \textbf{Fragmentação Horizontal (\textit{File Sharding / Byte Slicing}):} Divisão binária de cada arquivo enviado em duas metades (\texttt{.part1} e \texttt{.part2}), distribuídas entre os servidores centrais;
    \item \textbf{Balanceamento de Carga (\textit{Load Balancing}):} Distribuição uniforme de requisições utilizando o algoritmo \textbf{\textit{Round-Robin}} \textit{Thread-Safe};
    \item \textbf{Validação Distribuída de Duplicidade via Soma SHA-256:} Utilização de \textit{Worker Threads} nos servidores centrais para calcular o hash criptográfico SHA-256 de cada metade do arquivo e validar, pela soma composta dos hashes, se o arquivo completo já existe no sistema;
    \item \textbf{Exclusão Mútua Distribuída (\textit{Token Ring}):} Coordenação em anel lógico entre os servidores centrais para acesso exclusivo à \textbf{Região Crítica} (atualização do Banco de Metadados SQLite e replicação);
    \item \textbf{Replicação Primário-Backup e Tolerância a Falhas (\textit{Failover}):} Cópia síncrona das fatias para um nó de backup e recuperação automática e transparente durante o download caso um dos servidores centrais fique indisponível.
\end{itemize}

% =========================================================================
% 2. APRESENTAÇÃO DO PROBLEMA (JUSTIFICATIVA)
% =========================================================================
\section{Apresentação do Problema (Justificativa)}

Em um sistema distribuído de upload e download de arquivos acessado simultaneamente por múltiplos clientes, surgem cinco problemas clássicos que justificam a arquitetura proposta neste trabalho:

\subsection{Sobrecarga e Gargalo em Servidor Único}
Se todos os clientes enviarem arquivos grandes diretamente para um único servidor de armazenamento, a largura de banda, a CPU e o subsistema de disco desse nó rapidamente se tornarão um gargalo (\textit{bottleneck}). Faz-se necessário um nó intermediário de borda que distribua a carga e fragmente os arquivos entre múltiplos servidores de backend.

\subsection{Manutenção de Estado de Sessão em Ambiente Distribuído}
Sistemas que armazenam sessões de usuário na memória local de um servidor violam a escalabilidade horizontal, exigindo afinidade de sessão (\textit{sticky sessions}). O sistema precisa autenticar os usuários de forma \textbf{\textit{Stateless} (sem estado)}, permitindo que qualquer requisição seja processada de maneira independente.

\subsection{Condição de Corrida (\textit{Race Condition}) na Região Crítica}
Quando um arquivo é dividido ao meio e enviado simultaneamente em threads paralelas para o \texttt{central-1} e o \texttt{central-2}, ambos os servidores precisam registrar a localização de sua respectiva fatia no \textbf{Banco de Dados de Metadados} compartilhado (\texttt{metadata\_chunks.db} / \texttt{catalog.json}). Caso os dois servidores escrevam simultaneamente no banco de metadados sem coordenação, ocorrerá corrupção de dados, perda de registros ou bloqueio de escrita (\textit{database lock}). É indispensável um algoritmo de \textbf{Exclusão Mútua Distribuída} para serializar o acesso à Região Crítica de escrita, mantendo as leituras (downloads) livres e paralelas.

\subsection{Armazenamento Redundante de Arquivos Duplicados e Falsos Positivos em Fatias}
Usuários frequentemente fazem upload do mesmo arquivo repetidas vezes (inclusive alterando apenas o nome do arquivo). Entretanto, como o sistema divide os arquivos ao meio (50\% no Servidor 1 e 50\% no Servidor 2), validar apenas uma metade isolada geraria um \textbf{falso positivo}: dois arquivos diferentes que possuam os primeiros 50\% idênticos (ex.: mesmo cabeçalho de um PDF) e os outros 50\% diferentes não são arquivos duplicados. Portanto, o problema exige que os \textit{workers} de cada servidor central calculem o hash SHA-256 de sua metade e que a validação de duplicidade ocorra sobre a \textbf{combinação (soma) dos hashes das duas partes}.

\subsection{Perda de Dados por Queda de Nó (\textit{Single Point of Failure})}
Ao fragmentar um arquivo em duas partes distribuídas entre \texttt{central-1} e \texttt{central-2}, a falha ou desligamento de apenas um desses containers impediria a remontagem do arquivo original durante o download. O sistema precisa garantir \textbf{Tolerância a Falhas} através de um container de redundância (\texttt{server-backup}) e um mecanismo de \textbf{\textit{Failover} Automático}.

% =========================================================================
% 3. METODOLOGIA (COMO SERÁ RESOLVIDO O PROBLEMA)
% =========================================================================
\section{Metodologia}

Para solucionar os problemas delimitados na seção anterior, a metodologia adotada combina seis técnicas clássicas de Sistemas Distribuídos:

\begin{enumerate}
    \item \textbf{Autenticação \textit{Stateless} com JWT (HS256):}
    O cliente autentica-se no endpoint \texttt{POST /api/auth/login} do Servidor de Borda e recebe um token JWT assinado digitalmente contendo a identidade do usuário (\texttt{sub}) e seu tempo de expiração (\texttt{exp}). Em todas as requisições subsequentes de upload, listagem e download, o Servidor de Borda valida apenas a assinatura criptográfica do token no cabeçalho \texttt{Authorization: Bearer}, sem guardar estado de sessão em memória.

    \item \textbf{Tradução de Protocolo (REST $\rightarrow$ gRPC com Protocol Buffers):}
    Para oferecer compatibilidade universal aos clientes externos (navegadores, \textit{Swagger UI}, scripts HTTP), a borda utiliza \textbf{REST sobre HTTP/1.1}. Internamente, toda a comunicação entre o Servidor de Borda, os Servidores Centrais e o Servidor de Backup utiliza \textbf{gRPC sobre HTTP/2} com serialização binária \textbf{\textit{Protocol Buffers}}, configurada para suportar mensagens binárias de até 200 MB por fatia.

    \item \textbf{Fragmentação Binária (\textit{Byte Slicing}) e Balanceamento \textit{Round-Robin}:}
    Ao receber um arquivo de tamanho $N$ bytes, o Servidor de Borda calcula o ponto médio $\lfloor N / 2 \rfloor$ e fatia o fluxo de bytes em duas partes:
    \[
        \text{Parte 1} = \text{bytes}[0 : \lfloor N/2 \rfloor], \quad \text{Parte 2} = \text{bytes}[\lfloor N/2 \rfloor : N]
    \]
    O balanceador de carga \textit{Round-Robin} (protegido por \texttt{threading.Lock}) seleciona o servidor central de destino para cada fatia, alternando circularmente entre \texttt{central-1:50051} e \texttt{central-2:50052}.

    \item \textbf{Deduplicação Distribuída por Soma de Hashes SHA-256 nos Workers:}
    Antes de gravar os dados em disco, o Servidor de Borda aciona paralelamente o procedimento remoto \texttt{ComputeChunkHash} nos dois Servidores Centrais. Uma \textit{Worker Thread} do \texttt{central-1} calcula $H_1 = \text{SHA256}(\text{Parte 1})$ e uma \textit{Worker Thread} do \texttt{central-2} calcula $H_2 = \text{SHA256}(\text{Parte 2})$. A partir de ambos os resultados, calcula-se o \textbf{Hash Composto (Soma)}:
    \[
        H_{\text{composto}} = \text{SHA256}(H_1 \mathbin{\Vert} \text{``:''} \mathbin{\Vert} H_2)
    \]
    Dessa forma, se um arquivo novo tiver apenas 50\% de conteúdo igual a outro arquivo já existente, $H_{\text{composto}}$ será diferente e o upload será permitido. Se 100\% do arquivo (ambas as metades) for igual, $H_{\text{composto}}$ coincidirá com o registro existente no banco SQLite e o upload será bloqueado com status HTTP \texttt{409 Conflict}.

    \item \textbf{Exclusão Mútua Distribuída com Algoritmo \textit{Token Ring}:}
    Os servidores \texttt{central-1} e \texttt{central-2} formam um anel lógico no qual um \textbf{Token} circula continuamente via chamadas gRPC (\texttt{PassToken}) gerenciadas por uma thread dedicada em background.
    \begin{itemize}
        \item \textbf{Operações de Leitura (Download e Listagem):} Não alteram estado e são executadas imediatamente em paralelo, sem aguardar o Token;
        \item \textbf{Operações de Escrita (\texttt{SaveChunk}):} A thread trabalhadora fica aguardando em uma variável de condição (\texttt{threading.Condition}) até que o servidor detenha o Token. Ao obter o Token, a thread entra na \textbf{Região Crítica}, verifica a duplicidade de $H_{\text{composto}}$ no SQLite, grava a fatia em seu diretório exclusivo, replica para o \texttt{server-backup}, registra os ponteiros no banco SQLite e libera a Região Crítica para que o Token siga para o próximo nó do anel.
        \item \textbf{\textit{Self-Healing} no Anel:} Caso um dos servidores centrais caia, o nó sobrevivente detecta a falha na chamada \texttt{PassToken} e retém o Token localmente para não interromper o funcionamento do cluster.
    \end{itemize}

    \item \textbf{Replicação Primário-Backup e \textit{Failover} na Remontagem (\textit{Reassembly}):}
    Durante a execução na Região Crítica, cada Servidor Central invoca o método gRPC \texttt{ReplicateChunk} no container \texttt{server-backup:50053}, que armazena cópias redundantes tanto da \texttt{.part1} quanto da \texttt{.part2}. No momento do download, o Servidor de Borda consulta o Mapa de Chunks no banco SQLite e busca cada metade em seu servidor primário; caso o servidor primário esteja offline (\texttt{grpc.RpcError}), o Servidor de Borda redireciona automaticamente a busca da fatia faltante para o \texttt{server-backup} (\textit{Failover}), remonta os bytes em memória (\texttt{Parte 1 + Parte 2}) e entrega o arquivo íntegro ao usuário.
\end{enumerate}

% =========================================================================
% 4. ARQUITETURA DA SOLUÇÃO
% =========================================================================
\section{Arquitetura da Solução}

A arquitetura da solução é composta por quatro containers Docker interconectados por uma rede virtual do tipo \textit{bridge} (\texttt{sd\_network}). A Figura \ref{fig:arquitetura} apresenta o diagrama detalhado da arquitetura e o fluxo de comunicação entre os componentes.

\begin{figure}[htbp]
\centering
\begin{tikzpicture}[
    node distance=1.4cm and 1.8cm,
    box/.style={rectangle, draw=blue!60!black, fill=blue!5, thick, minimum width=3.6cm, minimum height=1.3cm, align=center, rounded corners=4pt, font=\footnotesize\sffamily},
    edgebox/.style={rectangle, draw=orange!80!black, fill=orange!10, thick, minimum width=4.2cm, minimum height=2.2cm, align=center, rounded corners=4pt, font=\footnotesize\sffamily},
    centralbox/.style={rectangle, draw=green!50!black, fill=green!8, thick, minimum width=4.2cm, minimum height=2.2cm, align=center, rounded corners=4pt, font=\footnotesize\sffamily},
    backupbox/.style={rectangle, draw=purple!70!black, fill=purple!8, thick, minimum width=4.2cm, minimum height=1.6cm, align=center, rounded corners=4pt, font=\footnotesize\sffamily},
    dbbox/.style={cylinder, shape border rotate=90, draw=black!70, fill=gray!15, thick, minimum width=3.2cm, minimum height=1.3cm, align=center, font=\scriptsize\sffamily},
    arrow/.style={-{Stealth[length=2.5mm]}, thick, draw=black!80},
    tokenarrow/.style={-{Stealth[length=2.5mm]}, ultra thick, draw=red!75!black, dashed},
    failoverarrow/.style={-{Stealth[length=2.5mm]}, thick, draw=purple!80!black, dotted}
]

% Cliente
\node[box] (client) {\textbf{Cliente Externo}\\(Swagger UI / REST Client)\\Token JWT Bearer};

% Edge Server
\node[edgebox, right=1.8cm of client] (edge) {\textbf{Container 1: edge-server}\\(FastAPI - Porta 8080)\\[2pt]
$\bullet$ Autenticação Stateless JWT\\
$\bullet$ Fatiamento Binário (50\% / 50\%)\\
$\bullet$ Load Balancer (Round-Robin)\\
$\bullet$ Soma SHA-256 \& Failover};

% Central 1 e Central 2
\node[centralbox, above right=0.6cm and 2.2cm of edge] (c1) {\textbf{Container 2: central-1}\\(gRPC Server - Porta 50051)\\[2pt]
$\bullet$ Pool 10 Worker Threads\\
$\bullet$ Cálculo SHA-256 (Parte 1)\\
$\bullet$ Thread Token Ring\\
$\bullet$ Pasta: \texttt{storage\_data/central-1}};

\node[centralbox, below right=0.6cm and 2.2cm of edge] (c2) {\textbf{Container 3: central-2}\\(gRPC Server - Porta 50052)\\[2pt]
$\bullet$ Pool 10 Worker Threads\\
$\bullet$ Cálculo SHA-256 (Parte 2)\\
$\bullet$ Thread Token Ring\\
$\bullet$ Pasta: \texttt{storage\_data/central-2}};

% Server Backup
\node[backupbox, below=1.6cm of edge] (backup) {\textbf{Container 4: server-backup}\\(gRPC Server - Porta 50053)\\[2pt]
$\bullet$ Réplica Redundante (\texttt{.part1} e \texttt{.part2})\\
$\bullet$ Pasta: \texttt{storage\_data/server-backup}};

% Banco de Metadados SQLite
\node[dbbox, right=1.6cm of $(c1.south)!0.5!(c2.north)$] (db) {\textbf{Região Crítica}\\Banco SQLite\\(\texttt{metadata\_chunks.db})\\Mapa de Ponteiros};

% Setas de comunicação
\draw[arrow] (client) -- node[above, font=\scriptsize, align=center] {HTTP REST\\(JWT Auth)} (edge);

\draw[arrow] (edge.east) -- node[above, sloped, font=\scriptsize] {gRPC: Parte 1 (50\%)} (c1.west);
\draw[arrow] (edge.east) -- node[below, sloped, font=\scriptsize] {gRPC: Parte 2 (50\%)} (c2.west);

% Token Ring entre Central 1 e Central 2
\draw[tokenarrow] (c1.south) to[bend left=20] node[right, font=\scriptsize\bfseries, text=red!75!black] {Token Ring (gRPC PassToken)} (c2.north);
\draw[tokenarrow] (c2.north) to[bend left=20] (c1.south);

% Acesso à Região Crítica (SQLite)
\draw[arrow] (c1.east) -- node[above, sloped, font=\scriptsize] {Escrita c/ Token} (db.north west);
\draw[arrow] (c2.east) -- node[below, sloped, font=\scriptsize] {Escrita c/ Token} (db.south west);

% Replicação para o Backup
\draw[arrow, draw=purple!70!black] (c1.south west) to[bend right=15] node[above, sloped, font=\scriptsize, text=purple!80!black] {gRPC ReplicateChunk} (backup.north east);
\draw[arrow, draw=purple!70!black] (c2.west) to[bend left=15] (backup.east);

% Failover do Edge para o Backup
\draw[failoverarrow] (edge.south) -- node[left, font=\scriptsize, text=purple!80!black, align=right] {gRPC GetChunk\\(Failover Automático)} (backup.north);

\end{tikzpicture}
\caption{Diagrama da Arquitetura Distribuída dos 4 Containers e Fluxo de Dados.}
\label{fig:arquitetura}
\end{figure}

\subsection{Descrição dos Componentes do Diagrama}

\begin{enumerate}
    \item \textbf{Container 1 --- Servidor de Borda (\texttt{edge-server} na porta \texttt{8080}):}
    Atua como \textit{API Gateway}, autenticador \textit{Stateless} JWT, fragmentador de arquivos (\textit{Sharding}) e balanceador de carga (\textit{Round-Robin}). Ele recebe o arquivo via REST, divide-o ao meio, solicita em paralelo aos workers dos servidores centrais o cálculo do SHA-256 de cada metade, gera o hash composto da soma e distribui as fatias via gRPC. No download, consulta o mapa de metadados, busca as fatias (acionando o \textit{Failover} no \texttt{server-backup} caso algum servidor central não responda) e remonta o arquivo original.

    \item \textbf{Container 2 e Container 3 --- Servidores Centrais (\texttt{central-1:50051} e \texttt{central-2:50052}):}
    São os nós primários de processamento e armazenamento distribuído. Cada container executa um servidor gRPC com um pool de \textbf{10 \textit{Worker Threads}} (\texttt{ThreadPoolExecutor}) para processar cálculos de hash SHA-256, leituras e escritas concorrentes, além de \textbf{1 \textit{Thread} dedicada em background} responsável pela circulação do \textbf{Token Ring}. Cada nó armazena apenas a sua metade correspondente do arquivo (\texttt{storage\_data/central-1/} guarda uma metade e \texttt{storage\_data/central-2/} guarda a outra).

    \item \textbf{Região Crítica --- Banco de Metadados (\texttt{metadata\_chunks.db} e \texttt{catalog.json}):}
    Banco de dados relacional SQLite que armazena apenas os metadados e ponteiros de localização dos arquivos e suas fatias (tabela \texttt{files} com o hash composto \texttt{combined\_file\_sha256} e tabela \texttt{file\_chunks} indicando em qual servidor central e em qual caminho de backup está cada parte). O acesso de escrita a este banco é protegido pelo algoritmo \textit{Token Ring}.

    \item \textbf{Container 4 --- Servidor de Redundância (\texttt{server-backup:50053}):}
    Container dedicado exclusivamente à alta disponibilidade e redundância de dados. Ele recebe cópias síncronas via gRPC (\texttt{ReplicateChunk}) das partes gravadas por \texttt{central-1} e \texttt{central-2}, mantendo-as em \texttt{storage\_data/server-backup/}. Quando um servidor central primário é desligado ou falha, o \texttt{server-backup} responde às chamadas \texttt{GetChunk} do Servidor de Borda, garantindo que nenhum arquivo fique inacessível.
\end{enumerate}

% =========================================================================
% 5. PROTÓTIPO DA SOLUÇÃO (DESCREVER O QUE FOI FEITO)
% =========================================================================
\section{Protótipo da Solução}

O protótipo foi inteiramente implementado em \textbf{Python 3.11} seguindo uma arquitetura limpa organizada em camadas (\texttt{config}, \texttt{domain}, \texttt{schemas}, \texttt{routes}, \texttt{controllers}, \texttt{services} e \texttt{infrastructure}) dentro de cada microsserviço.

\subsection{Organização em Camadas do Código-Fonte}

\begin{itemize}
    \item \textbf{\texttt{edge-server/}:}
    \begin{itemize}
        \item \texttt{src/routes/auth\_routes.py} e \texttt{file\_routes.py}: Definição das rotas HTTP REST e documentação interativa no \textit{Swagger UI} (\texttt{/docs});
        \item \texttt{src/services/auth\_service.py}: Geração e verificação de tokens JWT (\textit{PyJWT}) com esquema \texttt{HTTPBearer};
        \item \texttt{src/services/load\_balancer\_service.py}: Implementação do algoritmo \textit{Round-Robin} com exclusão mútua local (\texttt{threading.Lock});
        \item \texttt{src/services/edge\_gateway\_service.py}: Lógica de divisão binária do arquivo ao meio, disparo paralelo em \texttt{ThreadPoolExecutor} para cálculo dos hashes SHA-256 nos workers centrais, geração da soma composta dos hashes e lógica de \textit{Failover} automático no download;
        \item \texttt{src/infrastructure/grpc\_client.py}: Cliente gRPC configurado para transporte de blocos de até 200 MB.
    \end{itemize}

    \item \textbf{\texttt{central-server/}:}
    \begin{itemize}
        \item \texttt{proto/storage.proto}: Contrato formal \textit{Protocol Buffers} definindo os serviços \texttt{FileService} (\texttt{ComputeChunkHash}, \texttt{SaveChunk}, \texttt{ReplicateChunk}, \texttt{GetChunk}, \texttt{GetFileMap}, \texttt{ListFiles}) e \texttt{TokenRingService} (\texttt{PassToken});
        \item \texttt{src/services/token\_ring\_service.py}: Implementação do anel lógico com \texttt{threading.Condition}, controle de fila de escritores (\texttt{waiting\_writers}) e mecanismo \textit{Self-Healing} caso o nó vizinho fique offline;
        \item \texttt{src/services/file\_service.py}: Execução das \textit{Worker Threads} para cálculo do SHA-256 da fatia recebida, validação de duplicidade pela soma dos hashes dentro da Região Crítica e disparo da replicação para o backup;
        \item \texttt{src/infrastructure/storage\_repository.py}: Persistência física da fatia na pasta exclusiva do servidor e transações SQL no banco SQLite de metadados.
    \end{itemize}

    \item \textbf{\texttt{server-backup/}:}
    \begin{itemize}
        \item \texttt{src/services/backup\_service.py} e \texttt{src/infrastructure/backup\_repository.py}: Recebimento via gRPC das réplicas enviadas pelos servidores centrais e atendimento de leituras de emergência (\textit{Failover}).
    \end{itemize}
\end{itemize}

\subsection{Principais Trechos Implementados no Protótipo}

O Listing \ref{lst:sharding} demonstra o fatiamento binário do arquivo ao meio e o cálculo distribuído da soma dos hashes SHA-256 pelos workers dos servidores centrais no \texttt{edge-server}.

\begin{lstlisting}[language=Python, caption={Fatiamento ao meio e Soma Distribuída do SHA-256 (\texttt{edge\_gateway\_service.py}).}, label={lst:sharding}]
original_size = len(content)
midpoint = original_size // 2

# Divide os bytes do arquivo exatamente na metade
part1_bytes = content[:midpoint]
part2_bytes = content[midpoint:]

# Etapa 1: Workers do central-1 e central-2 calculam o SHA-256 de cada metade em paralelo
part1_hash = hash_results[1]["chunk_sha256"]
part2_hash = hash_results[2]["chunk_sha256"]

# Etapa 2: Soma dos Hashes (evita falso positivo caso apenas 50% do arquivo seja igual)
combined_raw = f"{part1_hash}:{part2_hash}".encode("utf-8")
combined_file_sha256 = hashlib.sha256(combined_raw).hexdigest()
\end{lstlisting}

O Listing \ref{lst:tokenring} exibe o controle de acesso à Região Crítica implementado no \texttt{TokenRingService} dos servidores centrais.

\begin{lstlisting}[language=Python, caption={Sincronização de Threads no Algoritmo Token Ring (\texttt{token\_ring\_service.py}).}, label={lst:tokenring}]
def enter_critical_section(self, operation_desc: str):
    with self.condition:
        self.waiting_writers += 1
        while not self.has_token or self.in_critical_section:
            self.condition.wait()
        self.waiting_writers -= 1
        self.in_critical_section = True

def exit_critical_section(self, operation_desc: str):
    with self.condition:
        self.in_critical_section = False
        self.condition.notify_all()
\end{lstlisting}

\subsection{Roteiro de Testes e Validação Realizada}

O protótipo foi validado através da interface \textit{Swagger UI} (\texttt{http://localhost:8080/docs}) e monitoramento em tempo real dos logs do Docker Compose:
\begin{enumerate}
    \item \textbf{Teste de Autenticação Stateless JWT:} Geração do token na rota \texttt{POST /api/auth/login} e autorização no cadeado do Swagger;
    \item \textbf{Teste de Upload, Sharding, Token Ring e Replicação:} Envio de um documento PDF (\texttt{Prolog (4).pdf}). Verificou-se no diretório \texttt{storage\_data/} que a primeira metade (\texttt{.part1}) foi gravada em \texttt{central-1/}, a segunda metade (\texttt{.part2}) foi gravada em \texttt{central-2/}, ambas foram replicadas em \texttt{server-backup/} e o banco \texttt{metadata\_chunks.db} registrou os ponteiros sem conflito de concorrência graças ao \textit{Token Ring};
    \item \textbf{Teste de Deduplicação pela Soma SHA-256:} Ao reenviar o mesmo arquivo, os workers calcularam os hashes das metades e o servidor central bloqueou a operação dentro da Região Crítica com status HTTP \texttt{409 Conflict}, identificando tratar-se de uma duplicata exata. Arquivos com apenas 50\% de conteúdo igual geraram uma soma SHA-256 distinta e foram aceitos normalmente;
    \item \textbf{Teste de Queda de Nó e \textit{Failover} Automático:} O container \texttt{central-2} foi interrompido (\texttt{docker stop central-2}) e solicitou-se o download do arquivo via \texttt{GET /api/download/\{filename\}}. O Servidor de Borda detectou a falha do \texttt{central-2}, recuperou a \texttt{Parte 1} do \texttt{central-1} e a \texttt{Parte 2} diretamente do \texttt{server-backup}, remontando o arquivo PDF perfeitamente e sem perda de dados.
\end{enumerate}

% =========================================================================
% 6. CONCLUSÃO
% =========================================================================
\section{Conclusão}

O desenvolvimento deste projeto permitiu aplicar, de forma prática e integrada, os principais conceitos teóricos da disciplina de Sistemas Distribuídos em um problema real de armazenamento, concorrência e disponibilidade de arquivos.

A separação da arquitetura em quatro containers Docker (\texttt{edge-server}, \texttt{central-1}, \texttt{central-2} e \texttt{server-backup}) demonstrou na prática as vantagens do desacoplamento e da transparência em sistemas distribuídos. Para o cliente externo, toda a complexidade interna fica oculta atrás de uma API REST simples e \textit{Stateless} autenticada via JWT. Internamente, o uso do \textbf{gRPC com \textit{Protocol Buffers}} proporcionou uma comunicação binária de alto desempenho tanto para o transporte das fatias de arquivos quanto para a passagem contínua do bastão no algoritmo \textbf{\textit{Token Ring}}.

Destaca-se ainda a eficácia de combinar três mecanismos complementares no backend: (1) a \textbf{fragmentação binária ao meio com balanceamento \textit{Round-Robin}}, que dividiu o custo de armazenamento e rede entre os dois servidores centrais; (2) a \textbf{validação de duplicidade pela soma dos hashes SHA-256 calculados pelos \textit{workers}}, que impediu a gravação de arquivos repetidos sem incorrer no erro de bloquear arquivos que possuam apenas uma metade semelhante; e (3) a \textbf{replicação síncrona para o servidor de backup aliada ao \textit{Failover} automático}, que comprovou a resiliência do sistema mesmo diante da queda completa de um dos servidores centrais.

\end{document}
